% Punto 2 · Introducción.
% Texto de la compañera, tomado de codigo/punto2/00_generar_html.py (tarjeta del sitio, apartado 2.1) y pasado a
% LaTeX tal cual; solo cambian el formato y los separadores (punto decimal, espacio para los miles).
% Lo que agrega la IA para unirlo va en \borradorIA.

\borradorIA{El punto 2 cambia de escala: de la cuenca del Fonce a una microcuenca periurbana de 1.0269~km$^2$ junto a
Denver (Colorado, EE.\,UU.), donde la lluvia que importa para el diseño es la de tormentas convectivas cortas. Con la lluvia horaria de una estación se construyen las curvas
intensidad--duración--frecuencia (IDF) y con ellas se dimensionan una alcantarilla (Problema A) y un tanque de
detención (Problema B).}

\paragraph{Sitio y datos.} Denver Stapleton. Estación: USW00023062. Variable: precipitación horaria. Unidades: mm.
Fuente: NOAA GHCNh, archivo histórico por estación. Las horas sin observación válida se mantienen como faltantes;
trazas se etiquetan aparte de los ceros medidos.

\pendiente{Contexto hidroclimático de Denver con bibliografía, motivación, objetivos y preguntas del punto 2. El
texto de la compañera no los trae.}
